% GENERATED FILE. Edit canonical lessons, then run npm run generate:books.
% canonical-source-hash: fnv1a64:8b7e05f671d3492e
% canonical-lessons: TA-C139-cirakam, TA-C139-kirai, TA-C139-palappalam, TA-C139-koyya, TA-C139-elumiccai, TA-R139-first-pass-food-and-nature, TA-R139-second-pass-nature-and-animals

\chapter{Cumin, Leafy greens, Jackfruit, Guava, Lemon}
\label{ch:ta-cumin-leafy-greens-jackfruit-guava-lemon}
\hlchaptermodality{139}

\begin{chapteropening}
\textbf{By the end of this chapter:} \emph{I can say cumin, leafy greens, a jackfruit, a guava and a lemon.}

Complete the last lesson of chapter 139: i can say cumin, leafy greens, a jackfruit, a guava and a lemon.
\end{chapteropening}

\section[cīrakam]{\ta{சீரகம்} (cīrakam) — cumin}

\label{lesson:TA-C139-cirakam}



\begin{quote}
Before the new one: say the Tamil for pepper, then the Tamil for mustard seed.
\end{quote}

\subsection*{You'll want to know: \ta{சீரகம்}}
\textbf{\ta{சீரகம்}} — \emph{cīrakam} — \textquotedblleft{}cumin\textquotedblright{}.

\begin{grammarlens}[title={What you've built}]
One more word for this chapter.
\end{grammarlens}

\subsection*{Your turn}
\begin{itemize}
\raggedright
  \item \emph{Say it:} \emph{cīrakam}
  \item \emph{Say it:} \emph{cīrakam}, once more
  \item \emph{Say it:} say \emph{cīrakam}
  \item \emph{Recall:} say the Tamil for pepper, then the Tamil for mustard seed, then say \emph{cīrakam} again
\end{itemize}

\begin{tcolorbox}[breakable,colback=teal!4,colframe=teal!35!black,title={Before you move on}]
What does \ta{சீரகம்} mean? (\textquotedblleft{}Cumin\textquotedblright{}.) Say it once more.
\end{tcolorbox}
\section[kīrai]{\ta{கீரை} (kīrai) — leafy greens}

\label{lesson:TA-C139-kirai}



\begin{quote}
Before the new one: say the Tamil for mustard seed, then the Tamil for cumin.
\end{quote}

\subsection*{You'll want to know: \ta{கீரை}}
\textbf{\ta{கீரை}} — \emph{kīrai} — \textquotedblleft{}leafy greens\textquotedblright{}.

\begin{grammarlens}[title={What you've built}]
One more word for this chapter.
\end{grammarlens}

\subsection*{Your turn}
\begin{itemize}
\raggedright
  \item \emph{Say it:} \emph{kīrai}
  \item \emph{Say it:} \emph{kīrai}, once more
  \item \emph{Say it:} say \emph{kīrai}
  \item \emph{Recall:} say the Tamil for mustard seed, then the Tamil for cumin, then say \emph{kīrai} again
\end{itemize}

\begin{tcolorbox}[breakable,colback=teal!4,colframe=teal!35!black,title={Before you move on}]
What does \ta{கீரை} mean? (\textquotedblleft{}Leafy greens\textquotedblright{}.) Say it once more.
\end{tcolorbox}
\section[palāppaḻam]{\ta{பலாப்பழம்} (palāppaḻam) — a jackfruit}

\label{lesson:TA-C139-palappalam}



\begin{quote}
Before the new one: say the Tamil for cumin, then the Tamil for leafy greens.

Say \textbf{\ta{எனக்குத்} \ta{தமிழ்} \ta{பிடிக்கும்}}, then the plain verb for wanting, \textbf{\ta{விரும்பு}}.
\end{quote}

\subsection*{You'll want to know: \ta{பலாப்பழம்}}
\textbf{\ta{பலாப்பழம்}} — \emph{palāppaḻam} — \textquotedblleft{}a jackfruit\textquotedblright{}.

\begin{grammarlens}[title={What you've built}]
One more word for this chapter.
\end{grammarlens}

\subsection*{Your turn}
\begin{itemize}
\raggedright
  \item \emph{Say it:} \emph{palāppaḻam}
  \item \emph{Say it:} \emph{palāppaḻam}, once more
  \item \emph{Say it:} say \emph{palāppaḻam}
  \item \emph{Recall:} say the Tamil for cumin, then the Tamil for leafy greens, then say \emph{palāppaḻam} again
\end{itemize}

\begin{tcolorbox}[breakable,colback=teal!4,colframe=teal!35!black,title={Before you move on}]
What does \ta{பலாப்பழம்} mean? (\textquotedblleft{}A jackfruit\textquotedblright{}.) Say it once more.
\end{tcolorbox}
\section[koyyā]{\ta{கொய்யா} (koyyā) — a guava}

\label{lesson:TA-C139-koyya}



\begin{quote}
Before the new one: say the Tamil for leafy greens, then the Tamil for a jackfruit.
\end{quote}

\subsection*{You'll want to know: \ta{கொய்யா}}
\textbf{\ta{கொய்யா}} — \emph{koyyā} — \textquotedblleft{}a guava\textquotedblright{}.

\begin{grammarlens}[title={What you've built}]
One more word for this chapter.
\end{grammarlens}

\subsection*{Your turn}
\begin{itemize}
\raggedright
  \item \emph{Say it:} \emph{koyyā}
  \item \emph{Say it:} \emph{koyyā}, once more
  \item \emph{Say it:} say \emph{koyyā}
  \item \emph{Recall:} say the Tamil for leafy greens, then the Tamil for a jackfruit, then say \emph{koyyā} again
\end{itemize}

\begin{tcolorbox}[breakable,colback=teal!4,colframe=teal!35!black,title={Before you move on}]
What does \ta{கொய்யா} mean? (\textquotedblleft{}A guava\textquotedblright{}.) Say it once more.
\end{tcolorbox}
\section[elumiccai]{\ta{எலுமிச்சை} (elumiccai) — a lemon}

\label{lesson:TA-C139-elumiccai}



\begin{quote}
Before the new one: say the Tamil for a jackfruit, then the Tamil for a guava.
\end{quote}

\subsection*{You'll want to know: \ta{எலுமிச்சை}}
\textbf{\ta{எலுமிச்சை}} — \emph{elumiccai} — \textquotedblleft{}a lemon\textquotedblright{}.

\begin{grammarlens}[title={What you've built}]
One more word for this chapter.
\end{grammarlens}

\subsection*{Your turn}
\begin{itemize}
\raggedright
  \item \emph{Say it:} \emph{elumiccai}
  \item \emph{Say it:} \emph{elumiccai}, once more
  \item \emph{Say it:} say \emph{elumiccai}
  \item \emph{Recall:} say the Tamil for a jackfruit, then the Tamil for a guava, then say \emph{elumiccai} again
\end{itemize}

\begin{tcolorbox}[breakable,colback=teal!4,colframe=teal!35!black,title={Before you move on}]
What does \ta{எலுமிச்சை} mean? (\textquotedblleft{}A lemon\textquotedblright{}.) Say it once more.
\end{tcolorbox}
\section[(dialogue)]{Food and nature — a first pass}

\label{lesson:TA-R139-first-pass-food-and-nature}



\begin{quote}
Say the words for a guava and a lemon.
\end{quote}

\subsection*{Your turn}
Say these aloud:

\begin{itemize}
\raggedright
  \item a parrot, a pigeon, an owl, a peacock, an ant
  \item a mosquito, the face, the mouth, the tongue, a finger
  \item a shoulder, the chest, a knee, hair, the skin
  \item a curtain, a mattress, a television, a mobile phone, a computer
\end{itemize}

\begin{tcolorbox}[breakable,colback=teal!4,colframe=teal!35!black,title={Before you move on}]
A parrot? (\textbf{\ta{கிளி}}, \emph{kiḷi}.) A pigeon? (\textbf{\ta{புறா}}, \emph{puṟā}.) An owl? (\textbf{\ta{ஆந்தை}}, \emph{āntai}.) A peacock? (\textbf{\ta{மயில்}}, \emph{mayil}.) An ant? (\textbf{\ta{எறும்பு}}, \emph{eṟumpu}.) A mosquito? (\textbf{\ta{கொசு}}, \emph{kocu}.) The face? (\textbf{\ta{முகம்}}, \emph{mukam}.) The mouth? (\textbf{\ta{வாய்}}, \emph{vāy}.) The tongue? (\textbf{\ta{நாக்கு}}, \emph{nākku}.) A finger? (\textbf{\ta{விரல்}}, \emph{viral}.) A shoulder? (\textbf{\ta{தோள்}}, \emph{tōḷ}.) The chest? (\textbf{\ta{மார்பு}}, \emph{mārpu}.) A knee? (\textbf{\ta{முழங்கால்}}, \emph{muḻaṅkāl}.) Hair? (\textbf{\ta{முடி}}, \emph{muṭi}.) The skin? (\textbf{\ta{தோல்}}, \emph{tōl}.) A curtain? (\textbf{\ta{திரைச்சீலை}}, \emph{tiraiccīlai}.) A mattress? (\textbf{\ta{மெத்தை}}, \emph{mettai}.) A television? (\textbf{\ta{தொலைக்காட்சி}}, \emph{tolaikkāṭci}.) A mobile phone? (\textbf{\ta{கைப்பேசி}}, \emph{kaippēci}.) A computer? (\textbf{\ta{கணினி}}, \emph{kaṇiṉi}.)
\end{tcolorbox}
\section[(dialogue)]{Nature and animals — a second pass}

\label{lesson:TA-R139-second-pass-nature-and-animals}



\begin{quote}
Say the words for an electric fan and scissors.
\end{quote}

\subsection*{Your turn}
Say these aloud:

\begin{itemize}
\raggedright
  \item an electric fan, scissors, a hammer, a nail, a ladder
  \item a bottle, a ball, a kite, a doll, a purse
  \item a ring, a bangle, a sock, a glove, a class
  \item an exam, an answer, tamarind, pepper, mustard seed
  \item cumin, leafy greens, a jackfruit, a guava, a lemon
\end{itemize}

\begin{tcolorbox}[breakable,colback=teal!4,colframe=teal!35!black,title={Before you move on}]
An electric fan? (\textbf{\ta{மின்விசிறி}}, \emph{miṉviciṟi}.) Scissors? (\textbf{\ta{கத்தரிக்கோல்}}, \emph{kattarikkōl}.) A hammer? (\textbf{\ta{சுத்தியல்}}, \emph{cuttiyal}.) A nail? (\textbf{\ta{ஆணி}}, \emph{āṇi}.) A ladder? (\textbf{\ta{ஏணி}}, \emph{ēṇi}.) A bottle? (\textbf{\ta{குப்பி}}, \emph{kuppi}.) A ball? (\textbf{\ta{பந்து}}, \emph{pantu}.) A kite? (\textbf{\ta{பட்டம்}}, \emph{paṭṭam}.) A doll? (\textbf{\ta{பொம்மை}}, \emph{pommai}.) A purse? (\textbf{\ta{பணப்பை}}, \emph{paṇappai}.) A ring? (\textbf{\ta{மோதிரம்}}, \emph{mōtiram}.) A bangle? (\textbf{\ta{வளையல்}}, \emph{vaḷaiyal}.) A sock? (\textbf{\ta{காலுறை}}, \emph{kāluṟai}.) A glove? (\textbf{\ta{கையுறை}}, \emph{kaiyuṟai}.) A class? (\textbf{\ta{வகுப்பு}}, \emph{vakuppu}.) An exam? (\textbf{\ta{தேர்வு}}, \emph{tērvu}.) An answer? (\textbf{\ta{பதில்}}, \emph{patil}.) Tamarind? (\textbf{\ta{புளி}}, \emph{puḷi}.) Pepper? (\textbf{\ta{மிளகு}}, \emph{miḷaku}.) Mustard seed? (\textbf{\ta{கடுகு}}, \emph{kaṭuku}.) Cumin? (\textbf{\ta{சீரகம்}}, \emph{cīrakam}.) Leafy greens? (\textbf{\ta{கீரை}}, \emph{kīrai}.) A jackfruit? (\textbf{\ta{பலாப்பழம்}}, \emph{palāppaḻam}.) A guava? (\textbf{\ta{கொய்யா}}, \emph{koyyā}.) A lemon? (\textbf{\ta{எலுமிச்சை}}, \emph{elumiccai}.)
\end{tcolorbox}
